\section{Resources and timing}
\begin{SDDTable}{lY}
Budget & Design or measured value\\\midrule
Generated Debug build & 31,452 bytes internal flash and 7,784 bytes RAM reported by linker on 16 October 2026.\\
Generated Release build & 17,012 bytes internal flash and 7,736 bytes RAM. Both exclude the unimplemented logger design.\\
Application workers & 10,240 bytes stacks total; USB adapter proposal adds 4096 bytes.\\
Four queue arrays & 11,264 bytes; queue objects, control events, mutex and status structures are additional.\\
USB chunks & Two 512-byte buffers: 1024 bytes plus descriptors.\\
Flash rate & 100 raw + 100 processed records/s; 80 bytes each, page padding: about 17,067 occupied bytes/s.\\
Ideal duration & About 245 s before metadata/session headers and other overhead. The board does not hold hours of logs at this rate.\\
Backlog & 256 queued records absorb about 1.28 s at 200 records/s. A queue cannot cure sustained overload.\\
\end{SDDTable}
The proposed stacks, queue arrays and chunk buffers total 26,624 bytes before middleware allocations and control objects. Do not simply add this to the generated RAM figure: pool replacement may remove or change existing allocations. Inspect the final map, stack high-water marks, middleware pools and interrupt/main stacks within 112 KiB.

The original generated port used a 100 Hz SysTick. The implementation now uses the shared 1 kHz configuration described below. Measure actual tick period and jitter and preserve the shared reload/constant across regeneration.

Use absolute 10 ms deadlines with wrap-safe comparisons. On overrun count it and skip to a defined future release rather than running an unlimited catch-up loop. At 8 MHz a 260-byte transaction takes 260 microseconds on the wire, excluding mutex and software overhead; measure actual worst-case lock duration. Flash internal busy waits release the bus.

Two program operations per record can consume 6 ms at the cited worst-case program time: 200 records/s can exceed sustained throughput before overhead. Qualification must justify actual rates, batching/commit changes, or a reduced logging rate. Pre-erase space in maintenance mode and preserve old sessions when full.

Implementation update (24 November 2026): \texttt{app\_timing.h} now supplies a shared 1000 Hz definition to the ThreadX configuration and SysTick assembly. Debug linkage reports approximately 97,520 bytes of RAM and 124 KiB of internal flash for the compact logger; measure stack high-water marks and timing on the board. Earlier generated-build figures above are historical baselines.

Compact logger update (24 November 2026): physical encoding 2 stores eight raw samples per 256-byte committed page, approximately 3200 bytes/s at 100 Hz. Processed records are reconstructed on export. A protected 20-second pre-trigger plus 300-second post-trigger capture occupies about 1 MiB. Ordinary 20-second segments roll automatically; protected events require explicit deletion. The earlier per-record throughput and duration figures describe the original design. The current output queue uses 16 KiB; final linker reports supersede the historical resource budgets above.
